\scanpage{181}\mspage{36}
pouvoir se paraphraser dans le cadre des schémas excellents, à grands coups de voisinages tubulaires étales. Dans l'un et l'autre cas, le ``type d'homotopie fin'' d'un espace modéré, resp. d'un schéma excellent, se définirait par passage à la limite à partir des ``\inserted{$\unsure{P|\mathcal{P}}$-}types d'homotopie'' associés à des stratifications équisingulières (à strates non sing.) $P$ de plus en plus fines. Ce type d'homotopie \inserted{``fin''} incarnerait la connaissance, non des seuls faisceaux ou \inserted{$n$-}champs loc constants, mais \inserted{(via un passage à la limite \unsure{inductive})} la connaissance de \uline{tous}. Et il dépendrait, en un sens convenable, fonctoriellement de $X$. Dans le cas d'un schéma de type fini sur un corps alg. clos disons, les \uline{théorèmes de finitude} cohomologique ou homotopique les plus forts s'exprimeraient précisément en termes des types d'homotopie fin, et diraient que les types d'homotopie \unsure{ordinaire} qui en sont les \unsure{élément|éléments} constituants sont essentiellement des ``\unsure{polyédres|polyèdres} finis'' -- et \unsure{$\overline{m}$|$=$} des variétés à bord compactes -- ou de façon plus précise,
% archive sheet number: 181
